
\subsection{Gröbner basis} 
The generating set $\ideal{LT(I)}$ was introduced in the last chapter. As we noted then, it is an large set. Dealing with monomials reduces the complexity drastically as opposed to polynomials, so if we can use the $\ideal{LT(I)}$ instead of the polynomial one, it would simplify calculations. Recall that the Hilbert Basis theorem stated that every polynomial ideal has a finite generating set. However, it is not unique. We will start with an ideal that is of special interest to the world of polynomial calculation.

\begin{definition}
    Choose a monomial ordering on the polynomial ring $\K[x]$. A finite nonzero subset, $G = \lbrace \vectorComponentsX{g}{t} \rbrace $ of an ideal, $I$, which is a subset of $\K[x]$ is a \textbf{Gröbner Basis} if 
    \begin{equation}
        \ideal{LT(g_1), ... , LT(g_t)} = \ideal{LT(I)}
    \end{equation}
\end{definition}

The construction is purposeful as now we can operate on a object of immense size ($\ideal{LT(I)}$) while restricting ourselvs to a finite number of generators. This makes the theoretical exploration less complex (due to the construction of $\ideal{LT(I)}$) while the instruction to a computer becomes guaranteed possible (as we are restricting algorithms input to finite). 

\subsubsection{Properties Of Gröbner Basis}
A problem that arises when performing multivariate division is that the order in which the polynomials appears affects the remainder of the quotient. This is particularly difficult when trying to establish deterministic outcomes. By fixing a lexographical order, gröbner basis gaurantee a unique remainder. 

\begin{proposition}
    Let $I \subseteq \polyRing$ be an ideal and let $G =  \lbrace \vectorComponentsX{g}{t} \rbrace$ be a Gröbner basus for I. Then given $ f \in \polyRing$, there is a unique $r \in  \polyRing$ with the two following properties:
    
    \begin{enumerate}
        \item No term of r is divisible by any of $LT(g_1), ... , LT(g_t)$.
        \item There is a $g \in I$ such that $f = g + r$.
    \end{enumerate}

    \begin{proof}
        Consider the division algorithm $f = q_1 g_1 \dots + q_n g_n + r $. Every $LT(g_t)$ is already divided out. In other words, if r had been divisible by the lead term there would have been a corrosponding $q_t$ to engulf it. That proves $(1)$. $(2)$ is solved by considering the fact that we are operating with gröbner basis; we must leverage our ability to show that $\specialGrobCondition$. To prove uniqueness, assume that the remainder is non-unique and use lemma \ref{lemma:4.2} to disprove that assumption. In other words: assume that $f = g + r = g' + r'$ follows our proposition. Rearranging, $r-r'= g-g'\in I$. Now, using the special property of the basis, $LT(r-r') \in \specialGrobCondition$. Since $LT(r-r')$ is on $LT(g_i)$, then it is also divisible by that same $LT(g_i)$. Using the first property, $r - r'$ cannot be divided by $LT(g_i)$, so in order for the division to succeed, $r - r' = 0$. This contridict the assumption that $r$ was non-unique.   
    \end{proof}
\end{proposition}

\subsubsection{Buchberger's Algorithm}
\input{subsections/GröbnerBases/BuchbergersAlgorithm.tex}